\documentclass[10pt]{article}
\usepackage[a4paper,margin=1in]{geometry}
\usepackage{graphicx}
\usepackage{booktabs}
\usepackage{pgfplotstable}
\usepackage{amsmath}
\usepackage{microtype}
\usepackage{hyperref}
\usepackage{newtxtext,newtxmath}
\pgfplotsset{compat=1.18}

\title{CEGVR: Candidate-First Verifiable Reasoning with Solver-Grounded Feedback}
\author{Stylianos Zacharioudakis}
\date{\today}

\begin{document}
\maketitle

\begin{abstract}
We present a linear-first, candidate-based LLM+SMT evaluation protocol for
verifiable reasoning under explicit feedback-granularity control. The main study
compares five methods: a one-shot reference, three compute-matched retry
baselines, and a new conflict-directed verifier-guided search method
(\texttt{cd\_vgs\_core\_rank}) that uses unsat-core structure for candidate
ranking and targeted repair. The implementation preserves the legacy trace-based
architecture as a compatibility and transfer layer while moving the main paper
path to a stricter candidate-verification interface. All experiments are
designed for local execution with an OpenAI-compatible endpoint such as
llama.cpp serving Qwen.
\end{abstract}

\section{Introduction}
The central question of this upgrade is not whether an LLM can emit a candidate,
but whether \emph{solver-grounded feedback} measurably improves verified
performance over mere repeated attempts under matched budgets. To answer that
question cleanly, the codebase now separates a candidate-first linear pipeline
from the preserved trace-based extension layer, standardizes outcome logging,
and reports paired arm comparisons with bootstrap confidence intervals.

\paragraph{Contributions.}
\begin{itemize}
  \item A candidate-first linear pipeline with explicit feedback-isolated retry baselines and a conflict-directed CoreRank flagship method.
  \item Paper-grade run records, metrics, and paired statistical comparisons.
  \item A preserved trace-based extension path, demonstrated via a small Sudoku appendix transfer experiment.
\end{itemize}

\section{Method}
\textbf{Main benchmark.} The primary benchmark is a solver-backed linear
feasibility dataset with SAT/UNSAT labels and structural difficulty metadata.

\textbf{Protocol.} The candidate generator returns either a full satisfying
assignment or an UNSAT claim. The verifier then either certifies the claim or
returns the richest solver-grounded diagnostic available. The main causal
comparison is among the four compute-matched multi-round methods; the one-shot
arm is retained as a repair-gain baseline.

\textbf{Local execution.} The documented local path uses Python 3.11, a local
OpenAI-compatible endpoint, and llama.cpp as the default serving stack. The
paper defaults are \(K=4\), \(R=4\), and thinking disabled.

\section{Results}
\subsection*{Results pending}
The complete planned model sweep has not been verified from artifacts in this
repository. No arm-level performance, statistical significance, transfer result
or method advantage is reported. Offline stub runs test implementation and are
not language-model experiments.

A future report must derive its tables from preserved per-problem records,
distinguish SAT assignment correctness from UNSAT classification, and report
errors, solver timeouts, actual tokens, latency and repair attempts. Round and
candidate caps alone do not demonstrate equal compute expenditure.


\section{Conclusion}
The upgraded repository now supports a linear-first, candidate-level
verifiable-reasoning study with strict feedback isolation and reproducible local
execution. The preserved trace layer remains useful for transfer evidence and
future extensions, but the main paper claims are now tied to the explicit
candidate-verification protocol.

\section*{Reproducibility}
The local paper pipeline is executed by \texttt{scripts/run\_all.sh}. It
bootstraps Python 3.11, regenerates the linear benchmark, runs the full arm
matrix, and rebuilds the figures, tables, and PDF.

\section*{Limitations}
The appendix transfer evidence is intentionally small and descriptive. The main
claims remain restricted to the linear benchmark and to the current local
serving setup. Additional training or larger transfer suites are outside the
scope of this slice.

\end{document}
